\pagenumbering{roman}
\begin{titlepage}
\thispagestyle{empty}
\begin{center}
  \vspace*{0.30\textheight}
  {\titlepagetitlefont MATLAB C-V Measurement Guideline\par}
  \vspace{0.45cm}
  {\titlepagesecondaryfont Zurich Instruments MFIA\par}
  \vspace{1.20cm}
  {\titlepagemetadatafont Author: R. Aliasgari Renani\par}
  \vspace{0.25cm}
  {\titlepagemetadatafont Date: 13 May 2026\par}
\end{center}
\end{titlepage}

\pagenumbering{arabic}
\setcounter{page}{1}

\section{Basic Use and Code Behavior}

\subsection{Purpose}

This document is a short lab reference for the MATLAB App Designer app
\path|CV_Zurich_3.mlapp|. The app connects to a Zurich Instruments MFIA through
LabOne, configures the MFIA impedance analyzer, sweeps the DC bias voltage, and
saves measured voltage and capacitance data to an Excel file.

The C-V measurement is done by applying a small AC impedance excitation at a
fixed frequency while sweeping a DC voltage across the device under test, for
example a diode or MOS capacitor. In the current code, the swept voltage is the
Signal Output offset and the saved capacitance is taken from
\path|/imps/<imps_channel>/sample| as \path|sample.param1|.

\subsection{Basic Use}

\begin{itemize}
  \item Start LabOne and make sure the MFIA is visible to the LabOne backend.
  \item The current app is hardcoded for device ID \path|dev5478|, API level 6,
  \path|127.0.0.1:8006|, and \path|PCIe| connection.
  \item Enter sample information, comments, save folder, impedance frequency,
  excitation amplitude, impedance model, voltage sweep limits, point count,
  settling, averaging, and bandwidth settings.
  \item Press \textbf{Start Experiment}. The app disables existing settings,
  configures the MFIA, runs the Sweeper, updates the plots, and then asks where
  to save the data.
  \item The Excel output has a \path|Settings| sheet and a \path|Measurements|
  sheet. The measurement columns are voltage and capacitance.
\end{itemize}

\subsection{Important Code Behavior}

\begin{itemize}
  \item The swept node is \path|sigouts/<out_c>/offset|. Therefore
  \path|start_voltage| and \path|stop_voltage| are DC output-offset values.
  \item The app reads the actual output channel from
  \path|/imps/<imps_channel>/output/select|. The MFIA node documentation lists
  this node as read-only and gives Signal Output 1 as the available output.
  \item The app hardcodes the generic signal output range and the impedance
  output range to 10 V.
  \item The app forces impedance output on at the end of configuration with
  \path|/imps/<imps_channel>/output/on = 1|.
  \item The app forces impedance input autorange on. The forced node/value is
  \path|/imps/<imps_channel>/auto/inputrange| = \path|1|.
  \item The app enables internal calibration and disables user compensation.
  \item The app sets impedance current input and voltage input to input 0.
  \item The app saves \path|sample.param1| as capacitance. This is only correct
  when the selected impedance model has capacitance as representation parameter
  2, such as an R-C representation.
\end{itemize}

\section{Parameter Reference}

\subsection{Signal Input Parameters}

\begin{itemize}
  \item \textbf{\path|in_c|}: Signal Input channel index used for
  \path|/sigins/<in_c>| and generic demodulator input selection. Default in the
  UI is 0. The exact available channel list should be read from the MFIA node
  tree. Source: app code and MFIA node docs.

  \item \textbf{\path|input_imp50|}: Sets
  \path|/sigins/<in_c>/imp50|. Allowed values are 0 for 10 MOhm and 1 for
  50 Ohm. UI default is 1. For high-impedance C-V samples, check whether 50 Ohm
  loading is really intended. Source: MFIA node docs.

  \item \textbf{\path|input_ac|}: Sets both
  \path|/sigins/<in_c>/ac| and \path|/imps/<imps_channel>/ac|. Allowed values
  are 0 for DC coupling and 1 for AC coupling. UI default is 1. AC coupling
  inserts a high-pass filter. Source: MFIA node docs.

  \item \textbf{\path|input_range|}: Sets
  \path|/sigins/<in_c>/range| in volts. UI default is 1 V. The manual says the
  range should exceed the incoming signal by roughly a factor of two and that
  the instrument selects the next higher available range. Exact available range
  values were not clearly stated in the checked sources.
\end{itemize}

\subsection{Signal Output Parameters}

\begin{itemize}
  \item \textbf{\path|out_c|}: Intended output channel. In the current code this
  is not really user-configurable. The app overwrites it with the read-only
  value from \path|/imps/<imps_channel>/output/select|. The checked MFIA node
  docs identify this as Signal Output 1.

  \item \textbf{\path|sigout_on|}: Sets \path|/sigouts/<out_c>/on| after setup.
  UI default is 1. Important: the impedance output is also forced on by the code,
  so this field does not fully control whether the MFIA drives the DUT.

  \item \textbf{\path|sigout_range|}: Misleading in the current app. The UI
  default is 10 V and the value is saved, but the code ignores the entered value
  and hardcodes \path|/sigouts/<out_c>/range = 10| and
  \path|/imps/<imps_channel>/output/range = 10|.

  \item \textbf{\path|sigout_amplitude|}: Sets the generic signal-output mixer
  amplitude at \path|/sigouts/<out_c>/amplitudes/<out_mixer_c>|. UI default is
  0.05 V. For this C-V app, \path|imps_amplitude| is the clearer impedance
  excitation amplitude.

  \item \textbf{\path|sigout_mixer_enable|}: Enables the selected generic
  signal-output mixer amplitude through
  \path|/sigouts/<out_c>/enables/<out_mixer_c>|. UI default is 1. This looks
  like a legacy lock-in style control and should be reviewed.
\end{itemize}

\subsection{Demodulator Parameters}

\begin{itemize}
  \item \textbf{\path|demod_phaseshift|}: Applies phase shift in degrees to
  \path|/demods/*/phaseshift|. UI default is 0. The checked sources did not
  clearly state a numeric range. Normally leave at 0 unless a calibrated phase
  correction is required.

  \item \textbf{\path|demod_order|}: Sets generic demodulator filter order and
  impedance demodulator order. Allowed values are 1 to 8, corresponding to
  6 dB/oct to 48 dB/oct. UI default is 4. Higher order affects settling time.

  \item \textbf{\path|demod_c|}: Generic demodulator index used for rate,
  harmonic, and enable. UI default is 0. Available demodulators depend on the
  MFIA options. Several other demodulator settings are applied to all
  demodulators, so this field is only a partial channel selector.

  \item \textbf{\path|demod_rate|}: Sets generic and impedance demodulator data
  rate in samples per second. UI default is 13000. The checked sources did not
  clearly state the numeric range, and Zurich notes that entered values may be
  approximated to supported values.

  \item \textbf{\path|demod_harmonic|}: Sets the reference harmonic for the
  generic and impedance demodulator. UI default is 1. Use 1 for ordinary
  fundamental C-V measurements unless there is a specific reason to measure a
  harmonic.

  \item \textbf{\path|demod_enable|}: Enables the selected generic demodulator.
  Allowed values are 0 off and 1 on. UI default is 1.

  \item \textbf{\path|demod_tc|}: Sets the generic and impedance demodulator
  time constant in seconds. UI default is 0.007 s. The checked sources did not
  clearly state the numeric range; the device exposes min/max time-constant
  properties in its node tree.
\end{itemize}

\subsection{Oscillator Parameters}

\begin{itemize}
  \item \textbf{\path|osc_c|}: Oscillator index used by the generic
  demodulators and the impedance demodulator. Documented oscillator selections
  are 0 to 3, depending on installed options. UI default is 0.

  \item \textbf{\path|osc_freq|}: Sets \path|/oscs/<osc_c>/freq|. UI default is
  3 MHz. The exact frequency range depends on the MFIA version/options. In this
  app, \path|imps_freq| is the more relevant C-V excitation frequency control.
\end{itemize}

\subsection{Impedance Analyzer Parameters}

\begin{itemize}
  \item \textbf{\path|imps_channel|}: Impedance analyzer unit index. UI default
  is 0. The manual describes one impedance analyzer unit normally, and two only
  with the MF-MD option.

  \item \textbf{\path|imps_amplitude|}: Sets
  \path|/imps/<imps_channel>/output/amplitude|. UI default is 0.01 V. This is
  the main AC test-signal amplitude for C-V. The checked node docs did not
  clearly state the exact numeric range.

  \item \textbf{\path|imps_freq|}: Sets \path|/imps/<imps_channel>/freq|. UI
  default is 100 kHz. It is the impedance excitation frequency for the sweep.
  The exact range depends on the MFIA model/options.

  \item \textbf{\path|imps_mode|}: Sets impedance measurement mode. Allowed
  values are 0 for 4-terminal and 1 for 2-terminal. UI default is 1. The manual
  notes that 2-terminal mode can be useful for very high impedance, while
  4-terminal mode removes series effects when wiring supports it.

  \item \textbf{\path|imps_model|}: Selects the impedance representation model.
  Important values for C-V are 0 for Rp parallel Cp and 1 for Rs + Cs. The full
  documented list is 0 to 10. The app saves \path|sample.param1|, so use a model
  where parameter 2 is capacitance.
\end{itemize}

\subsection{Sweep Parameters}

\begin{itemize}
  \item \textbf{\path|start_voltage|}: Sweeper start value for
  \path|sigouts/<out_c>/offset|. UI default is -0.5 V. This is the starting DC
  bias. The safe range depends on the MFIA output range, fixture, wiring, and
  DUT.

  \item \textbf{\path|stop_voltage|}: Sweeper stop value for
  \path|sigouts/<out_c>/offset|. UI default is 0.5 V. This is the final DC bias.

  \item \textbf{\path|sweep_mode_scan|}: Sweeper scan mode. Allowed values are
  0 sequential, 1 binary, 2 bidirectional, and 3 reverse. UI default is 2.
  Bidirectional mode is useful for checking hysteresis.

  \item \textbf{\path|sweep_samplecount|}: Number of sweep points. UI default is
  100. The checked Sweeper docs did not clearly state the maximum value.

  \item \textbf{\path|sweep_inaccuracy|}: Sweeper settling inaccuracy. Official
  Sweeper docs state the range as 1e-13 to 0.1. UI default is 0.001. Smaller
  values wait longer after each voltage step.

  \item \textbf{\path|xmapping|}: Sweep point spacing. Allowed values are
  0 linear and 1 logarithmic. UI default is 0. Use linear for normal voltage
  sweeps, especially if the sweep crosses 0 V.

  \item \textbf{\path|loopcount|}: Number of sweeps to perform. UI default is 1.
  More than one loop can test repeatability, but the current saved data format
  is not ideal for multiple loops.

  \item \textbf{\path|averaging_tc|}: Minimum averaging time in demodulator time
  constants per sweep point. UI default is 50. Larger values reduce noise but
  increase measurement time.

  \item \textbf{\path|averaging_sample|}: Minimum number of samples per sweep
  point. UI default is 0. Zurich recommends enough samples for reliable
  statistics; leaving this at 0 makes \path|averaging_tc| dominate.

  \item \textbf{\path|bandwidthcontrol|}: Sweeper bandwidth mode. Allowed values
  are 0 manual, 1 fixed, and 2 automatic. UI default is 2. Automatic bandwidth
  can override manually chosen demodulator bandwidth settings.

  \item \textbf{\path|bandwidthoverlap|}: Allows bandwidth overlap between
  neighboring sweep points. UI default is 0. The exact numeric enumeration was
  not clearly stated in the checked Sweeper docs. This is more important for
  frequency sweeps than for a DC-bias C-V sweep.
\end{itemize}

\subsection{Metadata and Saving}

\begin{itemize}
  \item \textbf{\path|SampleInfo|}: Free text saved in the settings sheet. Use
  it for sample ID, device type, contact, area, wafer, or location.

  \item \textbf{\path|Comments|}: Free text saved in the settings sheet. Use it
  for fixture, illumination, temperature, probe, and measurement notes.

  \item \textbf{\path|DefaultFolder|}: Save folder. It is restored from the
  MATLAB preference \path|MyApp/DataFolder|. If no preference exists, the app
  creates a \path|MyAppData| folder under the MATLAB user path.

  \item \textbf{\path|additionalComments|}: Optional extra comment added when
  using the \textbf{Save Data + Additional Comment} button after a measurement.

  \item \textbf{\path|device_id|}: Hardcoded private app property, currently
  \path|dev5478|. This is not user-facing and should probably become a device
  selection field later.
\end{itemize}

\section{TODO}

\subsection{App Cleanup}

\begin{itemize}
  \item Disable or remove \path|out_c| as a user input because the code reads it
  from a read-only MFIA node.
  \item Disable or remove \path|sigout_range| because the code hardcodes 10 V.
  \item Clarify \path|sigout_on| because the impedance output is forced on.
  \item Decide whether \path|sigout_amplitude| and
  \path|sigout_mixer_enable| are still needed for this C-V workflow.
  \item Keep \path|imps_freq| as the main excitation-frequency control and avoid
  confusion with \path|osc_freq|.
  \item Add UI validation for voltage limits, positive point count, valid
  channel indices, valid model selection, and logarithmic sweep compatibility.
  \item Replace numeric 0/1 and enum fields with dropdowns or checkboxes.
  \item Save more raw MFIA data later, not only voltage and capacitance:
  impedance, model parameter 1, model parameter 2, frequency, drive amplitude,
  and confidence indicators if available.
\end{itemize}

\subsection{DLTS Implementation Roadmap}

For DLTS, the MFIA should be used as the capacitance-transient meter. The
LakeShore controller handles temperature, and MATLAB should coordinate both
instruments.

The extra instrument used in many older DLTS setups is usually a boxcar averager
or DLTS correlator. It is not a second lock-in demodulator. Its job is to take
the capacitance transient after a filling pulse and calculate a rate-window
signal. We do not need that as separate hardware for the first version. We can
record the full transient with the MFIA and calculate the DLTS signal later in
MATLAB.

The MFIA already measures the small-signal capacitance. Zurich also shows DLTS
examples where the MFIA records capacitance transients and can generate the bias
pulse. An external pulse generator is only needed if the MFIA pulse is not good
enough for our required voltage, timing, edge speed, or isolation.

Minimum setup:
\begin{itemize}
  \item MFIA with LabOne API access from MATLAB.
  \item LakeShore temperature controller connected to MATLAB.
  \item Temperature stage or cryostat with stable wiring to the DUT.
  \item Low-parasitic fixture and compensation where needed.
  \item MATLAB code for transient averaging and rate-window calculation.
\end{itemize}

The current C-V app uses the LabOne Sweeper to step the DC output offset. DLTS
should be a separate mode. It should not sweep voltage point by point. It should
sit at one temperature, apply the same filling pulse many times, and record the
capacitance recovery after each pulse.

Inputs needed for DLTS:
\begin{itemize}
  \item reverse-bias voltage;
  \item filling-pulse voltage and pulse width;
  \item repetition period between pulses;
  \item transient length and DAQ sample rate;
  \item number of repeats to average;
  \item temperature start, stop, step, and stability tolerance;
  \item rate-window times for the DLTS signal.
\end{itemize}

One temperature point should run like this:
\begin{itemize}
  \item Set the LakeShore temperature and wait until it is stable.
  \item Configure the MFIA capacitance measurement: frequency, AC amplitude,
  model, bandwidth, and data rate.
  \item Apply reverse bias and wait for the baseline capacitance to settle.
  \item Arm LabOne DAQ for the capacitance sample node.
  \item Apply the filling pulse and return to reverse bias.
  \item Record the capacitance transient after the pulse.
  \item Repeat and average the transient.
  \item Save the raw transients as well as the averaged transient.
\end{itemize}

Processing in MATLAB:
\begin{itemize}
  \item Convert MFIA data to time, capacitance, temperature, and bias arrays.
  \item Reject bad traces: lost trigger, overload, clipping, contact jump, or
  unstable temperature.
  \item Average good traces at the same temperature.
  \item Calculate a simple DLTS signal first, for example
  \path|C(t1) - C(t2)| or two averaged time windows.
  \item Plot DLTS signal versus temperature for each rate window.
  \item Use the peak positions later for emission-rate and Arrhenius analysis.
  \item Leave Laplace-DLTS for later. It needs cleaner data and more careful
  numerical work.
\end{itemize}

First tests before real samples:
\begin{itemize}
  \item Check the pulse waveform before connecting a sensitive device.
  \item Test timing and triggering with a dummy capacitor or RC circuit.
  \item Repeat one temperature point many times and check drift/noise.
  \item Run a short temperature scan on a known diode or Schottky sample.
  \item Compare the capacitance level with ordinary C-V results.
\end{itemize}

For DLTS, saving only the final capacitance values is not enough. Each run
should save raw transients, averaged transients, temperature history, reverse
bias, pulse voltage, pulse width, repetition period, MFIA settings, DAQ settings,
and selected rate windows.

Practical implementation order:
\begin{enumerate}
  \item Add a separate \path|runDLTSExperiment| function.
  \item Reuse the MFIA connection/setup code.
  \item Replace the sweeper with triggered DAQ acquisition.
  \item Add the LakeShore temperature loop.
  \item Save raw transients first.
  \item Add simple two-point DLTS processing.
  \item Add plots for transient, averaged transient, and DLTS signal versus
  temperature.
\end{enumerate}
